% Originally: content/week-05.tex, \section{Lagrange multipliers} (Corsin Week 5)

\section{Second-order conditions on the constraint set}
\label{sec:second_order_lagrange}

% Generator: Claude Opus 5 (High) --- lead-in for the theorem below.
Every example so far has closed the same gap in the same way. The Lagrange condition returns a
list of candidates, and compactness plus a comparison of values decides among them. That works
only when the constraint set is compact, and it decides only the global extrema. Is there a
second-order test, of the kind \cref{thm:hessian_test} provides in the unconstrained case? There
is, and it is that same test, applied to the Lagrangian rather than to $f$, and read only along
directions that stay on the constraint set.

% Generator: Claude Sonnet 5 (High) --- new content: a second-order test for constrained
% extrema was never given anywhere in this document, only the unconstrained Hessian test
% (\cref{thm:hessian_test}) and the first-order Lagrange condition
% (\cref{prop:constrained_optimization}) that only locates \emph{candidates}. This fills that
% gap using exactly the two tools already available: the Lagrangian and the unconstrained
% Hessian test, applied after restricting to the constraint's tangent space.
% Feedback: Gemini 3.1 Pro (2026-08-11) --- named the single most important addition in this half
% of the document: the tutors' notes carry only the unconstrained Hessian test, so without this
% theorem the chapter can find critical points and not classify them. Keep it, and keep
% \cref{ex:second_order_test_on_circle}, which is what shows why restricting to the tangent space
% resolves a Hessian that is indefinite on all of $\mathbb{R}^n$.
\begin{theorem}[Second-order sufficient conditions for constrained extrema]
\label{thm:second_order_lagrange_test}
Let $U \subseteq \mathbb{R}^n$ be open, $f \in C^2(U,\mathbb{R})$, $g = (g_1,\dots,g_k) \in
C^2(U,\mathbb{R}^k)$, and suppose $x_0 \in g^{-1}\{0\}$ satisfies:
\begin{enumerate}[label=\textbf{(\alph*)}]
  \item $\nabla g_1(x_0),\dots,\nabla g_k(x_0)$ are linearly independent, and
  \item $\nabla f(x_0) = \sum_{j=1}^k\lambda_j\nabla g_j(x_0)$ for some $\lambda \in
    \mathbb{R}^k$ (a Lagrange critical point, as furnished by
    \cref{prop:constrained_optimization}).
\end{enumerate}
Write $L(x,\lambda) := f(x) - \sum_{j=1}^k\lambda_jg_j(x)$ for the Lagrangian, and let
\[T_{x_0} := \{v \in \mathbb{R}^n : Dg(x_0)v = 0\}\]
be the tangent space of the constraint set $g^{-1}\{0\}$ at $x_0$ (a $(n-k)$-dimensional
subspace, by \textbf{(a)}). Then:
\begin{enumerate}[label=\textbf{(\alph*)}]
  \item If $v\transp\Hess_xL(x_0,\lambda)v > 0$ for all $v \in T_{x_0}\setminus\{0\}$, then
    $x_0$ is a \textbf{strict local minimum} of $f|_{g^{-1}\{0\}}$.
  \item If $v\transp\Hess_xL(x_0,\lambda)v < 0$ for all $v \in T_{x_0}\setminus\{0\}$, then
    $x_0$ is a \textbf{strict local maximum} of $f|_{g^{-1}\{0\}}$.
  \item If $v\transp\Hess_xL(x_0,\lambda)v$ takes both signs on $T_{x_0}\setminus\{0\}$, then
    $x_0$ is \textbf{not} a local extremum of $f|_{g^{-1}\{0\}}$.
\end{enumerate}
Here $\Hess_xL(x_0,\lambda)$ denotes the Hessian of $L(\cdot,\lambda)$ in the $x$-variable
alone, evaluated at $x_0$.
\end{theorem}

\begin{proof}
By \textbf{(a)} and the implicit function theorem (\cref{thm:implicit_function_theorem}), after
relabelling coordinates so that the $k\times k$ block $\partial g/\partial(x_{n-k+1},\dots,x_n)$
is invertible at $x_0$, there is an open $V \subseteq \mathbb{R}^{n-k}$ and a $C^2$ map
$\varphi : V \to g^{-1}\{0\}$ with $\varphi(0) = x_0$, solving the constraint locally, such that
$D\varphi(0) : \mathbb{R}^{n-k} \to T_{x_0}$ is an isomorphism (this is the standard local
parametrization of the submanifold $g^{-1}\{0\}$; see also
\cref{thm:regular_value_theorem} for why it is one).

Since $g(\varphi(t)) \equiv 0$ on $V$, the terms $-\lambda_jg_j(\varphi(t))$ vanish identically,
so
\[h(t) := f(\varphi(t)) = L(\varphi(t),\lambda) \qquad \text{for all } t \in V.\]
By the chain rule, $Dh(0) = D_xL(x_0,\lambda)\,D\varphi(0) = 0$, since
$D_xL(x_0,\lambda) = \nabla f(x_0)\transp - \sum_j\lambda_j\nabla g_j(x_0)\transp = 0$ by
hypothesis \textbf{(b)}: so $h$ has a critical point at $0$, i.e.\ $f|_{g^{-1}\{0\}}$ has a
critical point at $x_0$, consistently with \textbf{(b)}.

Differentiating a second time, for $w \in \mathbb{R}^{n-k}$,
\[\Hess h(0)[w,w] = \big(D\varphi(0)w\big)\transp\Hess_xL(x_0,\lambda)\big(D\varphi(0)w\big)
  \ +\ D_xL(x_0,\lambda)\cdot D^2\varphi(0)[w,w],\]
and the second term vanishes because $D_xL(x_0,\lambda)=0$, regardless of what
$D^2\varphi(0)$ is. This is the same cancellation that already made $Dh(0)=0$. So, writing
$v := D\varphi(0)w$, which ranges over all of $T_{x_0}$ as $w$ ranges over $\mathbb{R}^{n-k}$
(since $D\varphi(0)$ is an isomorphism onto $T_{x_0}$),
\[\Hess h(0)[w,w] = v\transp\Hess_xL(x_0,\lambda)v.\]
The three cases now follow by reading the unconstrained Hessian test (\cref{thm:hessian_test})
off $h$ on the open set $V \subseteq \mathbb{R}^{n-k}$, and transporting the conclusion back
through $\varphi$, which is a local homeomorphism from $V$ onto a neighbourhood of $x_0$ in
$g^{-1}\{0\}$.
\begin{enumerate}[label=\textbf{(\alph*)}]
  \item $\Hess_xL(x_0,\lambda)$ positive definite on $T_{x_0}$ is exactly $\Hess h(0)$ positive
    definite, so the Hessian test gives a strict local minimum of $h$ at $0$, which transports
    to a strict local minimum of $f|_{g^{-1}\{0\}}$ at $x_0$.
  \item Identical, with every inequality reversed.
  \item If the quadratic form takes both signs on $T_{x_0}$, then $\Hess h(0)$ is indefinite, so
    $h$ increases along one direction through $0 \in V$ and decreases along another
    (\cref{thm:hessian_test} again), and transporting those two directions through $\varphi$
    shows $f|_{g^{-1}\{0\}}$ does the same near $x_0$.
\end{enumerate}
\end{proof}

\begin{remark}
This is exactly \cref{thm:hessian_test} run on the Lagrangian instead of on $f$ itself, with the
Hessian test's ordinary $\mathbb{R}^n$ replaced by the tangent space $T_{x_0}$. That is also why
$L$, not $f$, is the right function to differentiate: only $L$ has a genuine critical point at
$x_0$ (in the full, unrestricted sense), since $f$ itself typically does not. In practice
\textbf{(a)}--\textbf{(c)} are checked by picking any basis $v_1,\dots,v_{n-k}$ of $T_{x_0}$
(e.g.\ a basis of $\ker Dg(x_0)$) and testing the definiteness of the $(n-k)\times(n-k)$ matrix
$\big(v_i\transp\Hess_xL(x_0,\lambda)v_j\big)_{i,j}$ with \cref{thm:hessian_test}, or with the
determinant/trace recipe from \cref{sec:hessian_test}. There is no need to construct $\varphi$
explicitly, which was only a proof device.
\end{remark}

% Generator: Claude Sonnet 5 (High)
\begin{example}[The circle example, revisited with the second-order test]
\label{ex:second_order_test_on_circle}
Apply \cref{thm:second_order_lagrange_test} to \cref{ex:extrema_on_circle}
($f(x,y)=xy$ on $g(x,y):=x^2+y^2-2=0$) to see \emph{why} $(1,1)$ is a maximum and $(1,-1)$ a
minimum, without appealing to compactness. Here
\[\Hess_xL(x,y,\lambda) = \Hess f - \lambda\Hess g
  = \begin{pmatrix}0&1\\1&0\end{pmatrix} - \lambda\begin{pmatrix}2&0\\0&2\end{pmatrix}
  = \begin{pmatrix}-2\lambda&1\\1&-2\lambda\end{pmatrix},\]
which is indefinite (eigenvalues $-2\lambda\pm1$, of opposite sign whenever $|\lambda|<\tfrac12$)
or degenerate (at $\lambda=\pm\tfrac12$, exactly the values that occur here) on all of
$\mathbb{R}^2$, so the unconstrained Hessian test would say nothing. Restricted to the
$1$-dimensional tangent space it is a different story.

At $(1,1)$, $\lambda=\tfrac12$: $\Hess_xL = \left(\begin{smallmatrix}-1&1\\1&-1\end{smallmatrix}\right)$, and $T_{(1,1)} = \ker\begin{pmatrix}2&2\end{pmatrix} = \operatorname{span}\{(1,-1)\}$. For $v=(1,-1)$,
\[v\transp\Hess_xL\,v = \begin{pmatrix}1&-1\end{pmatrix}\begin{pmatrix}-1&1\\1&-1\end{pmatrix}\begin{pmatrix}1\\-1\end{pmatrix} = \begin{pmatrix}1&-1\end{pmatrix}\begin{pmatrix}-2\\2\end{pmatrix} = -4 < 0,\]
negative definite on $T_{(1,1)}$: a strict local maximum, matching $f(1,1)=1$.

At $(1,-1)$, $\lambda=-\tfrac12$: $\Hess_xL = \left(\begin{smallmatrix}1&1\\1&1\end{smallmatrix}\right)$, and $T_{(1,-1)} = \ker\begin{pmatrix}2&-2\end{pmatrix} = \operatorname{span}\{(1,1)\}$. For $v=(1,1)$,
\[v\transp\Hess_xL\,v = \begin{pmatrix}1&1\end{pmatrix}\begin{pmatrix}1&1\\1&1\end{pmatrix}\begin{pmatrix}1\\1\end{pmatrix} = \begin{pmatrix}1&1\end{pmatrix}\begin{pmatrix}2\\2\end{pmatrix} = 4 > 0,\]
positive definite on $T_{(1,-1)}$: a strict local minimum, matching $f(1,-1)=-1$. The same
computation at $(-1,-1)$ and $(-1,1)$ (by symmetry $f(-x,-y)=f(x,y)$) reproduces the same two
conclusions. This is the sharpest illustration available of the theorem's point: the
\emph{unrestricted} Hessian of $L$ is degenerate at both critical points, yet restricting to the
$1$-dimensional tangent space resolves it completely in both cases.
\end{example}

